\section{Nghiên cứu liên quan}
Các phương pháp truyền thống trong bài toán liên kết protein--phối tử chủ yếu dựa trên các thuật toán tìm kiếm không gian kết hợp với hàm chấm điểm để ước lượng năng lượng liên kết~\cite{dhillon2024molecular}. Mặc dù có ưu điểm về khả năng diễn giải dựa trên các nguyên lý vật lý, hóa học, các phương pháp này vẫn bộc lộ những hạn chế khi phải mô hình hóa sự biến đổi cấu dạng của protein, ảnh hưởng của dung môi và các tương tác phi tuyến phức tạp~\cite{li2024deep}.

Sự phát triển của học sâu đã mở ra nhiều hướng tiếp cận mới cho bài toán này. Một trong những phương pháp phổ biến là sử dụng mạng nơ-ron tích chập 3D~\cite{ragoza2017protein} trên biểu diễn dữ liệu dạng lưới voxel ba chiều~\cite{francoeur2020three, ragoza2017protein}. Các kiến trúc điển hình như GNINA~\cite{mcnutt2021gnina} cho thấy khả năng khai thác hiệu quả các đặc trưng hình học cục bộ và tương tác trong vùng lân cận. Tuy nhiên, các mô hình 3D CNN thường tiêu tốn nhiều tài nguyên bộ nhớ khi tăng độ phân giải và gặp khó khăn trong việc mô hình hóa các phụ thuộc không gian ở khoảng cách xa.

Một hướng tiếp cận khác là sử dụng mạng nơ-ron đồ thị~\cite{tang2025deep},
 trong đó cấu trúc phân tử được biểu diễn dưới dạng đồ thị với các đỉnh là nguyên tử và cạnh là các liên kết. Các mô hình như PotentialNet~\cite{feinberg2018potentialnet}, EquiBind~\cite{stark2022equibind} và TankBind~\cite{lu2022tankbind} đã cho thấy hiệu quả trong việc học biểu diễn cấu trúc và tương tác phân tử. Tuy nhiên, các kiến trúc đồ thị cơ bản thường thiếu khả năng khai thác trực tiếp ngữ cảnh không gian liên tục của vùng gắn kết nếu không được bổ sung các ràng buộc hình học chuyên biệt.

Gần đây, cơ chế chú ý và kiến trúc Transformer~\cite{vaswani2017attention} đã được sử dụng để giải quyết bài toán phụ thuộc xa giữa các biểu diễn đặc trưng, cho phép mô hình hoá các quan hệ không gian toàn cục một cách linh hoạt. Cùng với đó, học đa nhiệm~\cite{caruana1997multitask} cũng trở thành một tiếp cận phù hợp khi các nhiệm vụ có liên quan được tối ưu đồng thời, chẳng hạn phân loại tư thế và dự đoán ái lực. Việc chia sẻ biểu diễn giữa các nhiệm vụ có thể giúp mô hình khai thác các đặc trưng chung, liên kết các tri thức tốt hơn và tận dụng sự tương đồng giữa các nhiệm vụ.
